\begin{abstract}
Dissipation sets how long a quantum reservoir computer remembers its past inputs, and its rate is the usual control.
Here we show that its form matters as much, for how a memory is stored, written and read. When the qubits decay
independently, stronger loss shortens the memory. When they decay collectively, the quantum Zeno effect holds the
system in dark states that do not decay, and stronger loss prolongs the memory. A longer memory is harder to write,
because it leaves more time for new inputs to mix with the inputs already stored. We prove that the number of inputs
mixed into a stored signal is, on average, the memory lifetime times the fraction of the signal that one input mixes.
We also prove that this number stays fixed, however long the memory lasts, when the input amplitude falls as the
inverse square root of the lifetime. A weaker input leaves a fainter signal to read out, so the form of the
dissipation, the input amplitude and the measurement budget have to be chosen together.
\end{abstract}
